\RequirePackage{amsthm}
\documentclass[pdflatex,sn-basic]{sn-jnl}

\usepackage{lmodern}
\usepackage{amssymb}
\usepackage{mathtools}
\usepackage{array}
\usepackage{xcolor}
\usepackage{manyfoot}
\usepackage{booktabs}
\usepackage{enumitem}
\geometry{reset,left=2cm,right=2cm,top=2cm,bottom=2cm,bindingoffset=0cm}

\newcommand{\E}{\mathbb E}
\newcommand{\Pbb}{\mathbb P}
\newcommand{\1}{\mathbf 1}
\newcommand{\calF}{\mathcal F}
\newcommand{\calG}{\mathcal G}
\newcommand{\calH}{\mathcal H}
\newcommand{\dd}{\,\mathrm d}


\begin{document}

\title[Supplementary material for predictive HMM e-diagnostics]{Supplementary material for: Predictive e-diagnostics for multi-state movement models}

\author[1]{\fnm{Aurélien} \sur{Nicosia}}
\affil[1]{\orgname{Université Laval}, \orgaddress{\country{Canada}}}

\maketitle

\section{Simulation allocation}

The main article reports the scenarios that directly support the central methodological narrative: calibration, targeted movement diagnostics, predictable localization and predictable diagnostic combination. This supplement records additional scenarios that are useful for sensitivity analysis, implementation warnings or methodological extensions.

\begin{table}[h]
\centering
\footnotesize
\begin{tabular}{@{}>{\raggedright\arraybackslash}p{0.08\linewidth}>{\raggedright\arraybackslash}p{0.35\linewidth}>{\raggedright\arraybackslash}p{0.46\linewidth}@{}}
\toprule
Scenario & Placement & Reason \\
\midrule
S1 & Main article & Basic null calibration is required to support the e-process implementation. \\
S2 & Supplement & Parameter estimation uses an oracle-state estimator in the current simulation, so it is best treated as a sensitivity check. \\
S3 & Main article & The missing-state diagnostic is a core example of a full predictive HMM alternative. \\
S4 & Main article & Angular misspecification is a central movement-specific diagnostic. \\
S5 & Main article & Residual step-angle dependence is a central feature-level diagnostic. \\
S6 & Main article & Duration misspecification is a key movement failure and illustrates blockwise observable features. \\
S7 & Main article & Localization is a main theoretical and practical contribution. \\
S8 & Main article & Predictable mixtures and switching are main combination tools. \\
S9 & Supplement & The product example is an important warning, but it is not part of the recommended procedure. \\
S10 & Supplement & Long-horizon straightness is an advanced blockwise extension beyond the main simulation narrative. \\
S11 & Supplement & Individual cross-fitting is a validation-design extension. \\
S12 & Supplement & Composite envelopes are conservative and optional for the first article. \\
\bottomrule
\end{tabular}
\caption{Allocation of simulation scenarios between the main article and supplementary material.}
\label{tab:supp_allocation}
\end{table}

\section{Simulation DGP details}
\label{sec:dgp_details}

To ensure the simulation results can be understood independently, we briefly outline the data-generating processes (DGPs) used in the main scenarios. Von Mises distributions in the simulations are parametrized through circular standard deviations (sd); the corresponding concentration parameters $\kappa$ were obtained by numerical inversion of the circular variance formula $V = 1 - I_1(\kappa)/I_0(\kappa)$, where $I_p$ denotes the modified Bessel function of order $p$. All parameters are specified on the natural scale.

\begin{itemize}
    \item \textbf{S1 (Calibration)}: A 2-state HMM with initial probabilities $(0.5, 0.5)$ and transition matrix $\Gamma$ with $\Gamma(1,1)=0.8$, $\Gamma(2,2)=0.7$. State 1 (encamped) has step lengths drawn from $\text{Gamma}(\text{shape}=2, \text{rate}=0.5)$ and turning angles from $\text{von Mises}(\mu=0, \text{sd}=0.5)$. State 2 (exploratory) has step lengths $\text{Gamma}(10, 2)$ and turning angles $\text{von Mises}(\mu=0, \text{sd}=2)$. The diagnostic $q_t$ is an alternative 2-state HMM with slightly different parameters (e.g., transition and state-dependent parameters are perturbed relative to the null generator $p_0$).
    \item \textbf{S3 (Missing state)}: The true DGP is a 3-state HMM with an intermediate speed state, but the null model is restricted to $K=2$. True transition matrix $\Gamma$: row 1 $(0.91, 0.06, 0.03)$, row 2 $(0.08, 0.87, 0.05)$, row 3 $(0.08, 0.10, 0.82)$. Step shapes $(2.0, 7.5, 4.0)$, rates $(3.0, 2.0, 1.2)$, angle means $(0, 0, 0.6)$, angle SDs $(1.6, 0.35, 0.75)$. The diagnostic $q_t$ uses the \emph{true} 3-state DGP parameters \emph{(oracle diagnostic)}.
    \item \textbf{S4 (Angular misspecification)}: The true DGP has step shapes $(2.0, 7.0)$, rates $(3.0, 2.0)$, and angle SDs $(1.8, 0.25)$. The null model is misspecified with angle SDs $(0.85, 0.30)$ and state 2 centered at $\mu=0.25$ rather than $0$. The diagnostic $q_t$ is a fixed, manually chosen diffuse-angle alternative \emph{(fixed alternative)}.
    \item \textbf{S5 (Copula dependence)}: The true DGP includes a Gaussian copula with correlation $\rho = 0.6$ between step lengths and turning angles within each state, violating the conditional independence assumption of standard HMMs. The copula diagnostic uses a fixed target correlation $\rho_q = 0.55$ \emph{(fixed alternative)}.
    \item \textbf{S6 (Duration)}: The true DGP is a hidden semi-Markov model (HSMM) where the dwell time in both states follows a shifted Poisson distribution (mean 18), rather than a geometric distribution. The duration diagnostic is a blockwise observable feature (switch counts within each block) and is \emph{fixed} before seeing the data \emph{(fixed alternative)}.
    \item \textbf{S7 (Localized failure)}: Similar to S5 ($\rho = 0.8$ in state 2) but restricted to a time window (times 121 to 220). The diagnostic $q_t$ uses a fixed predictable weight concentrated on the failure window \emph{(fixed alternative)}.
    \item \textbf{S8 (Mixture switching)}: The alternative has two successive defects: angular misspecification (times 1 to 149) followed by step-angle copula dependence (times 150 to 300). The mixture and switching diagnostics combine fixed alternatives specified before the run \emph{(fixed alternatives)}.
\end{itemize}
Future work could further characterize diagnostic sensitivity by exploring weaker alternatives, such as scenarios with highly overlapping state-dependent distributions, only slightly misspecified turning angles (e.g., small differences in concentration parameters), moderate copula dependence ($\rho \approx 0.1$), or near-geometric dwell time distributions.

\section{Supplementary scenarios}

\subsection{S2: train/validation with estimated parameters}

Scenario S2 compares known parameters, oracle-state estimated parameters evaluated under the true generator and oracle-state estimated parameters evaluated under the fitted generator. At $\alpha=0.05$, the crossing rates are $0.057$, $0.053$ and $0.023$, respectively. These results are compatible with the split-validation interpretation, but the estimator uses simulated latent states. For that reason, S2 is not used as a main article result.

\begin{table}[h]
\centering
\footnotesize
\begin{tabular}{@{}lrrr@{}}
\toprule
Setting & Crossing rate & Monte Carlo SE & Mean final log e-value \\
\midrule
Known parameters, true generator & 0.057 & 0.013 & -23.868 \\
Oracle-state estimate, true generator & 0.053 & 0.013 & -23.885 \\
Oracle-state estimate, fitted generator & 0.023 & 0.009 & -24.565 \\
\bottomrule
\end{tabular}
\caption{S2 results at $\alpha=0.05$.}
\label{tab:supp_s2}
\end{table}

\subsection{S9: warning against naive parallel products}

Scenario S9 demonstrates the consequence of multiplying diagnostics computed in parallel on the same validation observations. At $\alpha=0.05$, a single valid e-process and the weighted average of two copies both cross at rate $0.030$, while naive products of two and three copies cross at rates $0.117$ and $0.215$. This supports the recommendation to use predictable mixtures, switching or weighted averages rather than parallel products.

\begin{table}[h]
\centering
\footnotesize
\begin{tabular}{@{}lrrr@{}}
\toprule
Method & Status & Crossing rate & Crossing rate minus $\alpha$ \\
\midrule
Single e-process & Valid & 0.030 & -0.020 \\
Weighted average of two copies & Valid & 0.030 & -0.020 \\
Naive product of two copies & Not valid by default & 0.117 & 0.067 \\
Naive product of three copies & Not valid by default & 0.215 & 0.165 \\
\bottomrule
\end{tabular}
\caption{S9 results at $\alpha=0.05$.}
\label{tab:supp_s9}
\end{table}

\subsection{S10: blockwise long-horizon straightness}

Scenario S10 uses block straightness, defined as net displacement divided by total path length, to detect a long-horizon failure generated by autocorrelated turning angles. At $\alpha=0.05$, the crossing rate is $0.036$ under the HMM null and $0.796$ under the long-horizon alternative. This scenario supports the blockwise theory but is kept in the supplement because it is a more advanced extension.

\begin{table}[h]
\centering
\footnotesize
\begin{tabular}{@{}lrrrr@{}}
\toprule
Generator & Crossing rate & Monte Carlo SE & Mean final log e-value & Mean straightness \\
\midrule
HMM null & 0.036 & 0.012 & -4.909 & 0.496 \\
Autocorrelated-angle alternative & 0.796 & 0.025 & 4.639 & 0.347 \\
\bottomrule
\end{tabular}
\caption{S10 results at $\alpha=0.05$.}
\label{tab:supp_s10}
\end{table}

\subsection{S11: individual validation and cross-fitted averages}

Scenario S11 studies validation by individuals. The weighted average of final individual e-values is conservative under the true and fitted null settings, with crossing rates $0.000$ and $0.004$ at $\alpha=0.05$. The exploratory scan over any individual is not a globally valid e-value and crosses at rates $0.304$ and $0.228$ under the two null settings. Under a single failed individual, the cross-fitted average crosses at rate $0.996$ and the individual scan crosses at rate $1.000$.

\begin{table}[h]
\centering
\footnotesize
\begin{tabular}{@{}llrr@{}}
\toprule
Scenario & Method & Status & Crossing rate \\
\midrule
True null & Cross-fitted final average & Valid final e-value & 0.000 \\
True null & Any individual scan & Exploratory & 0.304 \\
Fitted null & Cross-fitted final average & Valid final e-value & 0.004 \\
Fitted null & Any individual scan & Exploratory & 0.228 \\
Single failed individual & Cross-fitted final average & Valid final e-value & 0.996 \\
Single failed individual & Any individual scan & Exploratory & 1.000 \\
\bottomrule
\end{tabular}
\caption{S11 results at $\alpha=0.05$.}
\label{tab:supp_s11}
\end{table}

\subsection{S12: finite-family composite envelope}

Scenario S12 studies a conservative envelope over a finite family of three HMM null generators. It is important to note that the practical use of the composite envelope (Proposition 17) is primarily intended for finite or compactly restricted families. Continuous envelopes of unbounded parameter families for densities such as Gamma or von Mises can become numerically unstable or computationally prohibitive. The center-only denominator is not uniform over the family and gives a crossing rate of $0.980$ when the true generator is the diffuse-angle null. The composite envelope gives crossing rate $0.000$ at $\alpha=0.05$ for each generator in the finite family. The scenario illustrates uniform validity and its cost in conservatism.

\begin{table}[h]
\centering
\footnotesize
\begin{tabular}{@{}llrr@{}}
\toprule
Generator & Method & Status & Crossing rate \\
\midrule
Center null & Oracle true null & Valid for known generator & 0.020 \\
Center null & Center-only denominator & Not uniform over family & 0.020 \\
Center null & Composite envelope & Uniform finite-family envelope & 0.000 \\
Persistent null & Oracle true null & Valid for known generator & 0.027 \\
Persistent null & Center-only denominator & Not uniform over family & 0.000 \\
Persistent null & Composite envelope & Uniform finite-family envelope & 0.000 \\
Diffuse-angle null & Oracle true null & Valid for known generator & 0.043 \\
Diffuse-angle null & Center-only denominator & Not uniform over family & 0.980 \\
Diffuse-angle null & Composite envelope & Uniform finite-family envelope & 0.000 \\
\bottomrule
\end{tabular}
\caption{S12 results at $\alpha=0.05$.}
\label{tab:supp_s12}
\end{table}

\section{Additional proofs}
\label{sec:supp_proofs}

This section collects routine closure-property proofs that are used in the main article. The main text keeps the central density-ratio proof and the latent-state filtering argument, while the algebraic consequences are recorded here for completeness.

\subsection{Optional stopping and anytime p-values}

\begin{proof}[Proof of optional stopping and anytime-valid monitoring]
For the finite horizon $T<\infty$, Ville's maximal inequality applied to the nonnegative supermartingale $(E_{1:t})_{0\leq t\leq T}$ gives
\[
 \Pbb_0\left(\max_{0\leq t\leq T}E_{1:t}\geq \frac{1}{\alpha}\right)\leq \alpha.
\]
The logarithmic version follows by monotonicity of the logarithm. For any stopping time $\tau$, the stopped index $\tau\wedge T$ is bounded, so optional stopping for bounded nonnegative supermartingales gives
\[
 \E_0(E_{1:\tau\wedge T})\leq \E_0(E_{1:0})=1.
\]
Thus the stopped process is an e-value.
\end{proof}

\begin{proof}[Proof of the anytime p-value statement]
The event $\{\exists\,0\leq t\leq T:p_t^{\mathrm{any}}\leq\alpha\}$ is identical to
\[
 \left\{\sup_{0\leq t\leq T}E_{1:t}\geq\frac{1}{\alpha}\right\}.
\]
The result follows from the optional-stopping and anytime-valid monitoring bound.
\end{proof}

\subsection{Diagnostic mixtures, switching and sequential workflows}

\begin{proof}[Proof of predictable mixtures]
The mixture density
\[
q_t^{\mathrm{mix}}(y\mid\calF_{t-1})=\sum_{j=1}^Jw_{j,t}q_{j,t}(y\mid\calF_{t-1})
\]
is predictable. Since each component is a subdensity and the weights are predictable with $\sum_jw_{j,t}\leq1$,
\[
\int q_t^{\mathrm{mix}}(y\mid\calF_{t-1})\,\mu_t(\mathrm dy)
\leq \sum_{j=1}^Jw_{j,t}\leq1.
\]
The predictive e-process theorem therefore applies.
\end{proof}

\begin{proof}[Proof of predictable switching]
Set $w_{j,t}=\1\{J_t=j\}$. Since $J_t$ is $\calF_{t-1}$-measurable, these weights are predictable and sum to one. The result is the predictable-mixture result with a degenerate weight vector.
\end{proof}

\begin{proof}[Proof of sequential exploratory diagnostic workflows]
For each diagnostic stage $j$,
\[
 \E_0(D_{1:j}\mid\calG_{j-1})
 =D_{1:j-1}\E_0(D_j\mid\calG_{j-1})
 \leq D_{1:j-1}.
\]
Thus $(D_{1:m})_{0\leq m\leq M}$ is a nonnegative supermartingale with respect to $(\calG_m)$. Ville's inequality gives
\[
 \Pbb_0\left(\sup_{0\leq m\leq M}D_{1:m}\geq\frac{1}{\alpha}\right)\leq\alpha.
\]
\end{proof}

\begin{proof}[Proof of the weighted average and parallel-diagnostic bound]
Linearity of expectation gives
\[
\E_0E^{\mathrm{avg}}
=\sum_{j=1}^Jw_j\E_0E^{(j)}
\leq \sum_{j=1}^Jw_j\leq1.
\]
Thus $E^{\mathrm{avg}}$ is an e-value. If $\max_j w_jE^{(j)}\geq1/\alpha$, then $E^{\mathrm{avg}}\geq1/\alpha$, so Markov's inequality applied to $E^{\mathrm{avg}}$ gives the stated bound. The product warning follows because no conditional-independence or sequential-validity assumption has been imposed on the parallel factors.
\end{proof}

\subsection{Localization and tempering}

\begin{proof}[Proof of predictable weighted localization]
Since $E_t\geq0$ and $0\leq W_t\leq1$,
\[
E_t^{(W)}=(1-W_t)+W_tE_t\geq0.
\]
Moreover,
\[
\E_0(E_t^{(W)}\mid\calF_{t-1})
=1+W_t\{\E_0(E_t\mid\calF_{t-1})-1\}
\leq1.
\]
The cumulative product is therefore an e-process by the same supermartingale argument used for predictive e-processes.
\end{proof}

\begin{proof}[Proof of segment-specific localization]
Apply predictable weighted localization with $W_t=B_t\in\{0,1\}$.
\end{proof}

\begin{proof}[Proof of soft state-localized localization]
The filtered predictive probability $a_{t\mid t-1}(s)$ is $\calF_{t-1}$-measurable and lies in $[0,1]$, so it is an admissible predictable localization weight.
\end{proof}

\begin{proof}[Proof of power and linear tempering]
The linear case is predictable weighted localization with $W_t=\lambda_t$. For the power case, Jensen's inequality for the concave function $x\mapsto x^{\lambda_t}$ gives
\[
\E_0(E_t^{\lambda_t}\mid\calF_{t-1})
\leq \{\E_0(E_t\mid\calF_{t-1})\}^{\lambda_t}
\leq1.
\]
Both constructions therefore give conditional e-value increments.
\end{proof}

\subsection{Training, validation and composite nulls}

\begin{proof}[Proof of split-sample conditional validity]
Conditional on $\calH$, the fitted densities and diagnostic rules are fixed, and the validation filtration satisfies the assumptions of the predictive e-process theorem. The conditional anytime-valid bound follows from optional stopping. Taking expectations over $\calH$ gives the unconditional bound.
\end{proof}

\begin{proof}[Proof of the cross-fitted average]
For nonnegative weights satisfying $\sum_iw_i\leq1$,
\[
\E_0E^{\mathrm{cf}}
=\sum_{i=1}^nw_i\E_0E^{(i)}
\leq \sum_{i=1}^nw_i\leq1.
\]
No independence among the e-values is required for this averaging argument. A product would require additional conditional-independence or sequential-validity assumptions.
\end{proof}

\begin{proof}[Proof of the composite-null envelope]
For each fixed $\theta\in\Theta_0$, the envelope satisfies
\[
p_\theta(y\mid\calF_{t-1})\leq p_t^{\mathrm{env}}(y\mid\calF_{t-1})
\]
pointwise on the relevant support. Therefore,
\[
\E_\theta(\overline E_t\mid\calF_{t-1})
=\int
\frac{q_t(y\mid\calF_{t-1})}{p_t^{\mathrm{env}}(y\mid\calF_{t-1})}
p_\theta(y\mid\calF_{t-1})\,\mu_t(\mathrm dy)
\leq
\int q_t(y\mid\calF_{t-1})\,\mu_t(\mathrm dy)
\leq1.
\]
The cumulative product is therefore an e-process uniformly over the composite null family.
\end{proof}

\section{Reproducibility}

The simulation scripts are stored in the project directory \texttt{simulations/}. The complete simulation suite is run by \texttt{simulations/run\_all\_simulations.R}. Generated outputs are written to:
\begin{itemize}
\item \texttt{results/simulation\_tables/};
\item \texttt{results/simulation\_figures/}.
\end{itemize}
These generated outputs are intentionally ignored by Git; they should be regenerated from the scripts when preparing a submission archive.

\end{document}
